% sec2_1.tex — Section 2.1: Review of Limit Theorems for Sequences of Random Variables

The material of this section concerns the limiting behavior of a sequence of random
variables, $(z_1, z_2, \ldots)$.  Since the material may already be familiar to you, we
present it rather formally, in a series of definitions and theorems.  An authoritative
source is Rao (1973, Chapter 2c) which gives proofs of all the theorems included
in this section.  In this section and the rest of this book, a sequence $(z_1, z_2, \ldots)$
will be denoted by $\{z_n\}$.

%% ─── Various Modes of Convergence ──────────────────────────────────────────
\subsection*{Various Modes of Convergence}

\paragraph{Convergence in Probability.}
A sequence of random scalars $\{z_n\}$ \textbf{converges in probability} to a constant
(non-random) $\alpha$ if, for any $\varepsilon > 0$,
\begin{equation}
  \lim_{n \to \infty} \mathrm{Prob}(|z_n - \alpha| > \varepsilon) = 0.
  \label{eq:2.1.1}
\end{equation}
The constant $\alpha$ is called the \textbf{probability limit} of $z_n$ and is written as
``$\plim_{n \to \infty} z_n = \alpha$'' or ``$z_n \pto \alpha$''.  Evidently,
\[
  \text{``}z_n \pto \alpha\text{''} \quad \text{is the same as} \quad
  \text{``}z_n - \alpha \pto 0.\text{''}
\]

This definition of convergence in probability is extended to a sequence of random
vectors or random matrices (by viewing a matrix as a vector whose elements have
been rearranged) by requiring element-by-element convergence in probability.
That is, a sequence of $K$-dimensional random vectors $\{\mathbf{z}_n\}$ converges in
probability to a $K$-dimensional vector of constants $\bm{\alpha}$ if, for any $\varepsilon > 0$,
\begin{equation}
  \lim_{n \to \infty} \mathrm{Prob}(|z_{nk} - \alpha_k| > \varepsilon) = 0
  \quad \text{for all } k \ (= 1, 2, \ldots, K),
  \label{eq:2.1.2}
\end{equation}
where $z_{nk}$ is the $k$-th element of $\mathbf{z}_n$ and $\alpha_k$ the $k$-th element of $\bm{\alpha}$.

\paragraph{Almost Sure Convergence.}
A sequence of random scalars $\{z_n\}$ \textbf{converges almost surely} to a constant
$\alpha$ if
\begin{equation}
  \mathrm{Prob}\!\left(\lim_{n \to \infty} z_n = \alpha\right) = 1.
  \label{eq:2.1.3}
\end{equation}
We write this as ``$z_n \asto \alpha$.''  The extension to random vectors is analogous
to that for convergence in probability.  As will be mentioned below, this concept
of convergence is stronger than convergence in probability; that is, if a sequence
converges almost surely, then it converges in probability.  The concept involved
in \eqref{eq:2.1.3} is harder to grasp because the probability is about an event
concerning an infinite sequence $(z_1, z_2, \ldots)$.  For our purposes, however, all
that matters is that almost sure convergence is stronger than convergence in
probability.  If we can show that a sequence converges almost surely, that is one
way to prove the sequence converges in probability.

\paragraph{Convergence in Mean Square.}
A sequence of random scalars $\{z_n\}$ \textbf{converges in mean square} (or
\textbf{in quadratic mean}) to $\alpha$ (written as ``$z_n \to_{\text{m.s.}} \alpha$'') if
\begin{equation}
  \lim_{n \to \infty} \E[(z_n - \alpha)^2] = 0.
  \label{eq:2.1.4}
\end{equation}
The extension to random vectors is analogous to that for convergence in
probability: $\mathbf{z}_n \to_{\text{m.s.}} \bm{\alpha}$ if each element of $\mathbf{z}_n$
converges in mean square to the corresponding component of $\bm{\alpha}$.

\paragraph{Convergence to a Random Variable.}
In these definitions of convergence, the limit is a constant (i.e., a real number).
The limit can be a random variable.  We say that a sequence of $K$-dimensional
random variables $\{\mathbf{z}_n\}$ converges to a $K$-dimensional random variable
$\mathbf{z}$ and write $\mathbf{z}_n \pto \mathbf{z}$ if $\{\mathbf{z}_n - \mathbf{z}\}$
converges to $\mathbf{0}$:
\begin{align}
  \text{``}\mathbf{z}_n \pto \mathbf{z}\text{''} &\quad \text{is the same as} \quad
  \text{``}\mathbf{z}_n - \mathbf{z} \pto \mathbf{0}.\text{''} \label{eq:2.1.5a} \\
\intertext{Similarly,}
  \text{``}\mathbf{z}_n \asto \mathbf{z}\text{''} &\quad \text{is the same as} \quad
  \text{``}\mathbf{z}_n - \mathbf{z} \asto \mathbf{0},\text{''} \label{eq:2.1.5b} \\
  \text{``}\mathbf{z}_n \to_{\text{m.s.}} \mathbf{z}\text{''} &\quad \text{is the same as} \quad
  \text{``}\mathbf{z}_n - \mathbf{z} \to_{\text{m.s.}} \mathbf{0}.\text{''} \label{eq:2.1.5c}
\end{align}

\paragraph{Convergence in Distribution.}
Let $\{z_n\}$ be a sequence of random scalars and $F_n$ be the cumulative distribution
function (c.d.f.)\ of $z_n$.  We say that $\{z_n\}$ \textbf{converges in distribution}
to a random scalar $z$ if the c.d.f.\ $F_n$ of $z_n$ converges to the c.d.f.\ $F$ of
$z$ at every continuity point of $F$.\footnote{Do not be concerned about the
qualifier ``at every continuity point of $F$.''  For the most part, except possibly
for the chapters on the discrete or limited dependent variable, the relevant
distribution is continuous, and the distribution function is continuous at all points.}
We write ``$z_n \dto z$'' or ``$z_n \to_{\text{L}} z$'' and call $F$ the
\textbf{asymptotic} (or \textbf{limit} or \textbf{limiting}) \textbf{distribution}
of $z_n$.  Sometimes we write ``$z_n \dto F$,'' when the distribution $F$ is
well-known.  For example, ``$z_n \dto N(0,1)$'' should read ``$z_n \dto z$ and
the distribution of $z$ is $N(0,1)$ (normal distribution with mean 0 and variance 1).''
It can be shown from the definition that convergence in probability is stronger
than convergence in distribution, that is,
\begin{equation}
  \text{``}\mathbf{z}_n \pto \mathbf{z}\text{''} \;\Rightarrow\;
  \text{``}\mathbf{z}_n \dto \mathbf{z}.\text{''}
  \label{eq:2.1.6}
\end{equation}

A special case of convergence in distribution is that $z$ is a constant (a trivial
random variable).

The extension to a sequence of random vectors is immediate: $\mathbf{z}_n \dto \mathbf{z}$
if the joint c.d.f.\ $F_n$ of the random vector $\mathbf{z}_n$ converges to the joint
c.d.f.\ $F$ of $\mathbf{z}$ at every continuity point of $F$.  Note, however, that,
unlike the other concepts of convergence, for convergence in distribution,
element-by-element convergence does not necessarily mean convergence for the
vector sequence.  That is, ``each element of $\mathbf{z}_n \dto$ corresponding element
of $\mathbf{z}$'' does not necessarily imply ``$\mathbf{z}_n \dto \mathbf{z}$.''
A common way to establish the connection between scalar convergence and vector
convergence in distribution is

\medskip
\noindent\textbf{Multivariate Convergence in Distribution Theorem:}\quad
\emph{(stated in Rao, 1973, p.~128) Let $\{\mathbf{z}_n\}$ be a sequence of
$K$-dimensional random vectors.  Then:}
\[
  \text{``}\mathbf{z}_n \dto \mathbf{z}\text{''} \;\Leftrightarrow\;
  \text{``}\bm{\lambda}'\mathbf{z}_n \dto \bm{\lambda}'\mathbf{z}
  \text{ for any $K$-dimensional vector of real numbers.''}\!
\]

%% ─── Convergence in Distribution vs. Convergence in Moments ────────────────
\subsection*{Convergence in Distribution vs.\ Convergence in Moments}

It is worth emphasizing that the moments of the limit distribution of $z_n$ are not
necessarily equal to the limits of the moments of $z_n$.  For example,
``$z_n \dto z$'' does not necessarily imply
``$\lim_{n \to \infty} \E(z_n) = \E(z)$.''  However,

\begin{lemma}[convergence in distribution and in moments]
\label{lem:2.1}
Let $\alpha_{sn}$ be the $s$-th moment of $z_n$ and
$\lim_{n \to \infty} \alpha_{sn} = \alpha_s$ where $\alpha_s$ is finite
(i.e., a real number).  Then:
\[
  \text{``}z_n \dto z\text{''} \;\Rightarrow\;
  \text{``$\alpha_s$ is the $s$-th moment of $z$.''}\!
\]
\end{lemma}

Thus, for example, if the variance of a sequence of random variables converging
in distribution converges to some \emph{finite} number, then that number is the
variance of the limiting distribution.

%% ─── Relation among Modes of Convergence ───────────────────────────────────
\subsection*{Relation among Modes of Convergence}

Some modes of convergence are weaker than others.  The following theorem
establishes the relationship between the four modes of convergence.

\begin{lemma}[relationship among the four modes of convergence]
\label{lem:2.2}
\begin{enumerate}[label=(\alph*)]
  \item $\text{``}\mathbf{z}_n \to_{\emph{m.s.}} \bm{\alpha}\text{''} \;\Rightarrow\;
        \text{``}\mathbf{z}_n \pto \bm{\alpha}.\text{''}$
        So $\text{``}\mathbf{z}_n \to_{\emph{m.s.}} \mathbf{z}\text{''} \;\Rightarrow\;
        \text{``}\mathbf{z}_n \pto \mathbf{z}.\text{''}$
  \item $\text{``}\mathbf{z}_n \asto \bm{\alpha}\text{''} \;\Rightarrow\;
        \text{``}\mathbf{z}_n \pto \bm{\alpha}.\text{''}$
        So $\text{``}\mathbf{z}_n \asto \mathbf{z}\text{''} \;\Rightarrow\;
        \text{``}\mathbf{z}_n \pto \mathbf{z}.\text{''}$
  \item $\text{``}\mathbf{z}_n \pto \bm{\alpha}\text{''} \;\Leftrightarrow\;
        \text{``}\mathbf{z}_n \dto \bm{\alpha}.\text{''}$
        \emph{(That is, if the limiting random variable is a constant
        [a trivial random variable], convergence in distribution is the same
        as convergence in probability.)}
\end{enumerate}
\end{lemma}

%% ─── Three Useful Results ──────────────────────────────────────────────────
\subsection*{Three Useful Results}

Having defined the modes of convergence, we can state three results essential for
developing large-sample theory.

\begin{lemma}[preservation of convergence for continuous transformation]
\label{lem:2.3}
Suppose $\mathbf{a}(\cdot)$ is a vector-valued continuous function that does not
depend on $n$.
\begin{enumerate}[label=(\alph*)]
  \item $\text{``}\mathbf{z}_n \pto \bm{\alpha}\text{''} \;\Rightarrow\;
        \text{``}\mathbf{a}(\mathbf{z}_n) \pto \mathbf{a}(\bm{\alpha}).\text{''}$
        Stated differently,
        \[
          \plim_{n \to \infty} \mathbf{a}(\mathbf{z}_n)
          = \mathbf{a}\!\left(\plim_{n \to \infty} \mathbf{z}_n\right)
        \]
        provided the plim exists.
  \item $\text{``}\mathbf{z}_n \dto \mathbf{z}\text{''} \;\Rightarrow\;
        \text{``}\mathbf{a}(\mathbf{z}_n) \dto \mathbf{a}(\mathbf{z}).\text{''}$
\end{enumerate}
\end{lemma}

An immediate implication of Lemma~\ref{lem:2.3}(a) is that the usual arithmetic
operations preserve convergence in probability.  For example:
\begin{align*}
  &\text{``}x_n \pto \beta,\; y_n \pto \gamma\text{''} \;\Rightarrow\;
   \text{``}x_n + y_n \pto \beta + \gamma\text{''} \\
  &\text{``}x_n \pto \beta,\; y_n \pto \gamma\text{''} \;\Rightarrow\;
   \text{``}x_n y_n \pto \beta\gamma.\text{''} \\
  &\text{``}x_n \pto \beta,\; y_n \pto \gamma\text{''} \;\Rightarrow\;
   \text{``}x_n/y_n \pto \beta/\gamma,\text{''  provided that } \gamma \neq 0. \\
  &\text{``}\mathbf{Y}_n \pto \bm{\Gamma}\text{''} \;\Rightarrow\;
   \text{``}\mathbf{Y}_n^{-1} \pto \bm{\Gamma}^{-1},\text{'' provided that
   $\bm{\Gamma}$ is invertible.}
\end{align*}

The next result about combinations of convergence in probability and in
distribution will be used repeatedly to derive the asymptotic distribution of the
estimator.

\begin{lemma}
\label{lem:2.4}
\begin{enumerate}[label=(\alph*)]
  \item $\text{``}\mathbf{x}_n \dto \mathbf{x},\;
         \mathbf{y}_n \pto \bm{\alpha}\text{''} \;\Rightarrow\;
         \text{``}\mathbf{x}_n + \mathbf{y}_n \dto \mathbf{x} + \bm{\alpha}.\text{''}$
  \item $\text{``}\mathbf{x}_n \dto \mathbf{x},\;
         \mathbf{y}_n \pto \mathbf{0}\text{''} \;\Rightarrow\;
         \text{``}\mathbf{y}_n'\mathbf{x}_n \pto \mathbf{0}.\text{''}$
  \item $\text{``}\mathbf{x}_n \dto \mathbf{x},\;
         \mathbf{A}_n \pto \mathbf{A}\text{''} \;\Rightarrow\;
         \text{``}\mathbf{A}_n\mathbf{x}_n \dto \mathbf{A}\mathbf{x},\text{''}$
         provided that $\mathbf{A}_n$ and $\mathbf{x}_n$ are conformable.
         In particular, if\/ $\mathbf{x} \sim N(\mathbf{0}, \bm{\Sigma})$, then
         $\mathbf{A}_n\mathbf{x}_n \dto N(\mathbf{0},
         \mathbf{A}\bm{\Sigma}\mathbf{A}')$.
  \item $\text{``}\mathbf{x}_n \dto \mathbf{x},\;
         \mathbf{A}_n \pto \mathbf{A}\text{''} \;\Rightarrow\;
         \text{``}\mathbf{x}_n'\mathbf{A}_n^{-1}\mathbf{x}_n \dto
         \mathbf{x}'\mathbf{A}^{-1}\mathbf{x},\text{''}$
         provided that $\mathbf{A}_n$ and $\mathbf{x}_n$ are conformable and
         $\mathbf{A}$ is nonsingular.
\end{enumerate}
\end{lemma}

Parts (a) and (c) are sometimes called \textbf{Slutzky's Theorem}.  By setting
$\bm{\alpha} = \mathbf{0}$, part (a) implies:
\begin{equation}
  \text{``}\mathbf{x}_n \dto \mathbf{x},\;
  \mathbf{y}_n \pto \mathbf{0}\text{''} \;\Rightarrow\;
  \text{``}\mathbf{x}_n + \mathbf{y}_n \dto \mathbf{x}.\text{''}
  \label{eq:2.1.7}
\end{equation}
That is, if $\mathbf{z}_n = \mathbf{x}_n + \mathbf{y}_n$ and
$\mathbf{y}_n \pto \mathbf{0}$ (i.e., if $\mathbf{z}_n - \mathbf{x}_n \pto \mathbf{0}$),
then the asymptotic distribution of $\mathbf{z}_n$ is the same as that of
$\mathbf{x}_n$.  When $\mathbf{z}_n - \mathbf{x}_n \pto \mathbf{0}$, we sometimes
(but not always) say that the two sequences are \textbf{asymptotically equivalent}
and write it as
\[
  \text{``}\mathbf{z}_n \underset{a}{\sim} \mathbf{x}_n\text{''} \quad \text{or} \quad
  \text{``}\mathbf{z}_n = \mathbf{x}_n + o_p,\text{''}
\]
where $o_p$ is some suitable random variable ($\mathbf{y}_n$ here) that converges
to zero in probability.

A standard trick in deriving the asymptotic distribution of a sequence of random
variables is to find an asymptotically equivalent sequence whose asymptotic
distribution is easier to derive.  In particular, by replacing $\mathbf{y}_n$ by
$\mathbf{y}_n - \bm{\alpha}$ in part (b) of the lemma, we obtain
\begin{equation}
  \text{``}\mathbf{x}_n \dto \mathbf{x},\;
  \mathbf{y}_n \pto \bm{\alpha}\text{''} \;\Rightarrow\;
  \text{``}\mathbf{y}_n'\mathbf{x}_n \underset{a}{\sim}
  \bm{\alpha}'\mathbf{x}_n\text{''} \quad \text{or} \quad
  \text{``}\mathbf{y}_n'\mathbf{x}_n = \bm{\alpha}'\mathbf{x}_n + o_p.\text{''}
  \label{eq:2.1.8}
\end{equation}
The $o_p$ here is $(\mathbf{y}_n - \bm{\alpha})'\mathbf{x}_n$.  Therefore,
replacing $\mathbf{y}_n$ by its probability limit does not change the asymptotic
distribution of $\mathbf{y}_n'\mathbf{x}_n$, \emph{provided} $\mathbf{x}_n$
converges in distribution to some random variable.

The third result will allow us to test nonlinear hypotheses given the asymptotic
distribution of the estimator.

\begin{lemma}[the ``delta method'']
\label{lem:2.5}
Suppose $\{\mathbf{x}_n\}$ is a sequence of $K$-dimensional random vectors such
that $\mathbf{x}_n \pto \bm{\beta}$ and
\[
  \sqrt{n}(\mathbf{x}_n - \bm{\beta}) \dto \mathbf{z},
\]
and suppose $\mathbf{a}(\cdot)\colon \mathbb{R}^K \to \mathbb{R}^r$ has continuous
first derivatives with $\mathbf{A}(\bm{\beta})$ denoting the $r \times K$ matrix of
first derivatives evaluated at $\bm{\beta}$:
\[
  \underset{(r \times K)}{\mathbf{A}(\bm{\beta})}
  = \frac{\partial \mathbf{a}(\bm{\beta})}{\partial \bm{\beta}'}.
\]
Then
\[
  \sqrt{n}[\mathbf{a}(\mathbf{x}_n) - \mathbf{a}(\bm{\beta})]
  \dto \mathbf{A}(\bm{\beta})\mathbf{z}.
\]
In particular:
\[
  \text{``}\sqrt{n}(\mathbf{x}_n - \bm{\beta}) \dto N(\mathbf{0}, \bm{\Sigma})\text{''}
  \;\Rightarrow\;
  \text{``}\sqrt{n}[\mathbf{a}(\mathbf{x}_n) - \mathbf{a}(\bm{\beta})]
  \dto N\!\big(\mathbf{0},\,
  \mathbf{A}(\bm{\beta})\,\bm{\Sigma}\,\mathbf{A}(\bm{\beta})'\big).\text{''}
\]
\end{lemma}

\begin{proof}
By the mean-value theorem from calculus (see Section~7.3 for a statement of the
theorem), there exists a $K$-dimensional vector $\mathbf{y}_n$ between
$\mathbf{x}_n$ and $\bm{\beta}$ such that
\[
  \mathbf{a}(\mathbf{x}_n) - \mathbf{a}(\bm{\beta})
  = \underset{(r \times K)}{\mathbf{A}(\mathbf{y}_n)}
    \underset{(K \times 1)}{(\mathbf{x}_n - \bm{\beta})}.
\]
Multiplying both sides by $\sqrt{n}$, we obtain
\[
  \sqrt{n}[\mathbf{a}(\mathbf{x}_n) - \mathbf{a}(\bm{\beta})]
  = \mathbf{A}(\mathbf{y}_n)\sqrt{n}(\mathbf{x}_n - \bm{\beta}).
\]
Since $\mathbf{y}_n$ is between $\mathbf{x}_n$ and $\bm{\beta}$ and since
$\mathbf{x}_n \pto \bm{\beta}$, we know that $\mathbf{y}_n \pto \bm{\beta}$.
Moreover, the first derivative $\mathbf{A}(\cdot)$ is continuous by assumption.
So by Lemma~\ref{lem:2.3}(a),
\[
  \mathbf{A}(\mathbf{y}_n) \pto \mathbf{A}(\bm{\beta}).
\]
By Lemma~\ref{lem:2.4}(c), this and the hypothesis that
$\sqrt{n}(\mathbf{x}_n - \bm{\beta}) \dto \mathbf{z}$ imply that
\[
  \mathbf{A}(\mathbf{y}_n)\sqrt{n}(\mathbf{x}_n - \bm{\beta})
  \;\big(= \sqrt{n}[\mathbf{a}(\mathbf{x}_n) - \mathbf{a}(\bm{\beta})]\big)
  \dto \mathbf{A}(\bm{\beta})\mathbf{z}. \qedhere
\]
\end{proof}

%% ─── Viewing Estimators as Sequences of Random Variables ───────────────────
\subsection*{Viewing Estimators as Sequences of Random Variables}

Let $\hat{\bm{\theta}}_n$ be an estimator of a parameter vector $\bm{\theta}$ based
on a sample of size $n$.  The sequence $\{\hat{\bm{\theta}}_n\}$ is an example of a
sequence of random variables, so the concepts introduced in this section for
sequences of random variables are applicable to $\{\hat{\bm{\theta}}_n\}$.  We say
that an estimator $\hat{\bm{\theta}}_n$ is \textbf{consistent for} $\bm{\theta}$ if
\[
  \plim_{n \to \infty} \hat{\bm{\theta}}_n = \bm{\theta} \quad \text{or} \quad
  \hat{\bm{\theta}}_n \pto \bm{\theta}.
\]

The \textbf{asymptotic bias} of $\hat{\bm{\theta}}_n$ is defined as
$\plim_{n \to \infty} \hat{\bm{\theta}}_n - \bm{\theta}$.\footnote{Some authors use
the term ``asymptotic bias'' differently.  Amemiya (1985), for example, defines it
to mean $\lim_{n \to \infty} \E(\hat{\bm{\theta}}_n) - \bm{\theta}$.}
So if the estimator is consistent, its asymptotic bias is zero.  A consistent
estimator $\hat{\bm{\theta}}_n$ is \textbf{asymptotically normal} if
\[
  \sqrt{n}(\hat{\bm{\theta}}_n - \bm{\theta}) \dto N(\mathbf{0}, \bm{\Sigma}).
\]
Such an estimator is called $\bm{\sqrt{n}}$\textbf{-consistent}.  The acronym
sometimes used for ``consistent and asymptotically normal'' is \textbf{CAN}.  The
variance matrix $\bm{\Sigma}$ is called the \textbf{asymptotic variance} and is
denoted $\avar(\hat{\bm{\theta}}_n)$.  Some authors use the notation
$\avar(\hat{\bm{\theta}}_n)$ to mean $\bm{\Sigma}/n$ (which is zero in the limit).
In this book, $\avar(\hat{\bm{\theta}}_n)$ is the variance of the limiting
distribution of $\sqrt{n}(\hat{\bm{\theta}}_n - \bm{\theta})$.

%% ─── Laws of Large Numbers and Central Limit Theorems ──────────────────────
\subsection*{Laws of Large Numbers and Central Limit Theorems}

For a sequence of random scalars $\{z_i\}$, the \textbf{sample mean}
$\bar{z}_n$ is defined as
\[
  \bar{z}_n \equiv \frac{1}{n} \sum_{i=1}^{n} z_i.
\]
Consider the sequence $\{\bar{z}_n\}$.  Laws of large numbers (LLNs) concern
conditions under which $\{\bar{z}_n\}$ converges either in probability or almost
surely.  An LLN is called strong if the convergence is almost surely and weak if
the convergence is in probability.  We can derive the following weak LLN easily
from Part~(a) of Lemma~\ref{lem:2.2}.

\medskip
\noindent\textbf{A Version of Chebychev's Weak LLN:}
\[
  \text{``}\lim_{n \to \infty} \E(\bar{z}_n) = \mu,\;
  \lim_{n \to \infty} \Var(\bar{z}_n) = 0\text{''} \;\Rightarrow\;
  \text{``}\bar{z}_n \pto \mu.\text{''}
\]

\medskip
This holds because, under the condition specified, it is easy to prove (see an
analytical question) that $\bar{z}_n \to_{\text{m.s.}} \mu$.  The following strong
LLN assumes that $\{z_i\}$ is i.i.d.\ (independently and identically distributed),
but the variance does not need to be finite.

\medskip
\noindent\textbf{Kolmogorov's Second Strong Law of Large Numbers:}\quad
\emph{Let $\{z_i\}$ be i.i.d.\ with $\E(z_i) = \mu$.\footnote{So the mean exists
and is finite (a real number).  When a moment (e.g., the mean) is indicated, as
here, then by implication the moment is assumed to exist and is finite.}
Then $\bar{z}_n \asto \mu$.}

\medskip
These LLNs extend readily to random vectors by requiring element-by-element
convergence.

Central Limit Theorems (CLTs) are about the limiting behavior of the difference
between $\bar{z}_n$ and $\E(\bar{z}_n)$ (which equals $\E(z_i)$ if $\{z_i\}$ is
i.i.d.)\ blown up by $\sqrt{n}$.  The only Central Limit Theorem we need for the
case of i.i.d.\ sequences is:

\medskip
\noindent\textbf{Lindeberg--Levy CLT:}\quad
\emph{Let $\{\mathbf{z}_i\}$ be i.i.d.\ with
$\E(\mathbf{z}_i) = \bm{\mu}$ and $\Var(\mathbf{z}_i) = \bm{\Sigma}$.  Then}
\[
  \sqrt{n}(\bar{\mathbf{z}}_n - \bm{\mu})
  = \frac{1}{\sqrt{n}} \sum_{i=1}^{n} (\mathbf{z}_i - \bm{\mu})
  \dto N(\mathbf{0},\, \bm{\Sigma}).
\]

This reads: a sequence of random vectors
$\{\sqrt{n}(\bar{\mathbf{z}}_n - \bm{\mu})\}$ converges in distribution to a random
vector whose distribution is $N(\mathbf{0}, \bm{\Sigma})$.  (Usually, the
Lindeberg--Levy CLT is for a sequence of \emph{scalar} random variables.  The
vector version displayed above is derived from the scalar version as follows.
Let $\{\mathbf{z}_i\}$ be i.i.d.\ with $\E(\mathbf{z}_i) = \bm{\mu}$ and
$\Var(\mathbf{z}_i) = \bm{\Sigma}$, and let $\bm{\lambda}$ be any vector of real
numbers of the same dimension.  Then $\{\bm{\lambda}'\mathbf{z}_n\}$ is a sequence
of scalar random variables with $\E(\bm{\lambda}'\mathbf{z}_n) = \bm{\lambda}'\bm{\mu}$
and $\Var(\bm{\lambda}'\mathbf{z}_n) = \bm{\lambda}'\bm{\Sigma}\bm{\lambda}$.  The
scalar version of Lindeberg--Levy then implies that
\[
  \sqrt{n}(\bm{\lambda}'\bar{\mathbf{z}}_n - \bm{\lambda}'\bm{\mu})
  = \bm{\lambda}'\sqrt{n}(\bar{\mathbf{z}}_n - \bm{\mu})
  \dto N(0,\, \bm{\lambda}'\bm{\Sigma}\bm{\lambda}).
\]
But this limit distribution is the distribution of $\bm{\lambda}'\mathbf{x}$ where
$\mathbf{x} \sim N(\mathbf{0}, \bm{\Sigma})$.  So by the Multivariate Convergence
in Distribution Theorem stated a few pages back,
$\{\sqrt{n}(\bar{\mathbf{z}}_n - \bm{\mu})\} \dto \mathbf{x}$, which is the claim
of the vector version of Lindeberg--Levy.)
